\chapter{Teneurs de marché automatisés}\label{ch:m3:automated-market-makers}

Un pool détient de l'ether et des dollars et n'affiche aucun prix. Chacun peut lui envoyer un \omterm{def:m3:blockchains-for-traders:key}{jeton} et en retirer l'autre,
en quantité quelconque, à condition qu'après la transaction le produit de ses deux soldes ne soit pas inférieur à
celui d'avant. Cette seule règle, écrite en quelques dizaines de lignes d'un \omterm{def:m3:blockchains-for-traders:contract}{smart contract}, fait un marché~: le prix est
le rapport des soldes~; un acheteur le fait monter en retirant de l'ether, et les arbitragistes le maintiennent aligné
avec le prix sur les \omterm{def:m3:centralised-exchanges:cex}{plateformes centralisées} (CEX) en négociant dès que les deux diffèrent. Ceux qui déposent les
soldes perçoivent les frais et supportent un coût moins évident que ces frais et, pour la plupart des pools, plus élevé.
Ce chapitre présente le \omterm{def:m3:automated-market-makers:amm}{teneur de marché automatisé}~: la règle du produit constant et sa variante à \omterm{def:m3:automated-market-makers:concentrated}{liquidité concentrée},
l'invariant utilisé pour les \omterm{def:m3:blockchains-for-traders:stable}{stablecoins}, ce que gagnent et perdent les fournisseurs de liquidité, mesuré
selon la méthode de la littérature récente, ainsi que les routeurs et les enchères qui se placent entre les traders et les pools.

\section{AMM à produit constant}

\begin{definition}[Plateforme décentralisée, teneur de marché automatisé, pool de liquidité]\label{def:m3:automated-market-makers:amm}
Une \emph{plateforme décentralisée}\index{plateforme décentralisée} (decentralised exchange, DEX) est une plateforme de négociation mise en œuvre
sous forme de \omterm{def:m3:blockchains-for-traders:contract}{smart contracts} sur une \omterm{def:m3:blockchains-for-traders:chain}{blockchain}, sur laquelle les utilisateurs négocient depuis leur propre portefeuille, et chaque transaction se règle sur la chaîne.
Un \emph{teneur de marché automatisé}\index{teneur de marché automatisé} (automated market maker, AMM) est une plateforme décentralisée dont les prix sont
fixés par une formule appliquée aux soldes d'un \emph{pool de liquidité}\index{pool de liquidité}~: un contrat qui détient des
réserves de deux \omterm{def:m3:blockchains-for-traders:key}{jetons} ou plus, déposées par des fournisseurs de liquidité qui reçoivent en contrepartie un droit sur celles-ci et sur les frais.
\end{definition}

\begin{definition}[AMM à produit constant]\label{def:m3:automated-market-makers:cp}
Un \emph{AMM à produit constant}\index{AMM à produit constant} (constant-product market maker) est un \omterm{def:m3:automated-market-makers:amm}{teneur de marché automatisé} dont le
pool de réserves $x$ et $y$ n'accepte une transaction que si $xy$ ne diminue pas, avec des frais $\gamma_{\mathrm{fee}}$ prélevés sur
l'entrée et laissés dans le pool.
\end{definition}

\begin{proposition}[Sortie d'un swap et prix marginal]\label{prop:m3:automated-market-makers:swap}
Un trader qui envoie $\Delta x$ vers un \omterm{def:m3:automated-market-makers:cp}{AMM à produit constant} de réserves $(x, y)$ et de frais $\gamma_{\mathrm{fee}}$ reçoit
\[
\Delta y = \frac{y\,(1 - \gamma_{\mathrm{fee}})\,\Delta x}{x + (1 - \gamma_{\mathrm{fee}})\,\Delta x}.
\]
Le prix marginal du premier \omterm{def:m3:blockchains-for-traders:key}{jeton} dans le pool vaut $y/x$~; le prix moyen de la transaction, $\Delta y / \Delta x$, est inférieur
à $y/x$ en raison des frais et d'un impact sur le prix qui croît avec $\Delta x / x$. Sans frais, le produit est exactement conservé~;
avec frais, il augmente.
\end{proposition}

\begin{proof}
Seule la part $(1 - \gamma_{\mathrm{fee}})\Delta x$ compte pour l'invariant~: $(x + (1 - \gamma_{\mathrm{fee}})\Delta x)(y - \Delta y) = xy$
donne $\Delta y$. La part des frais de $\Delta x$ s'ajoute aux réserves par-dessus, donc $(x + \Delta x)(y - \Delta y) > xy$.
Le prix marginal vaut $-dy/dx$ le long de $xy = k$, soit $y/x$.
\end{proof}

Un pool de 1\,000 ether et 3~millions de dollars affiche l'ether à 3\,000~dollars. En vendant 10 ether avec des frais de 0.30~\%, le pool renvoie
29\,614.74~dollars, soit un prix moyen de 2\,961.47~: 128 points de base sous 3\,000, dont 30 pour les frais et environ 98 pour l'impact. Le pool n'a ni cours acheteur ni
cours vendeur~; sa « fourchette » vaut le double des frais plus l'impact de la taille, et sa profondeur est fixée par ses réserves. Son prix ne bouge
que lorsque quelqu'un y négocie. Quand l'ether bouge sur les \omterm{def:m3:centralised-exchanges:cex}{plateformes centralisées}, le prix du pool reste périmé jusqu'à
ce qu'un arbitragiste le remette en ligne en y négociant, et le profit de l'arbitragiste est prélevé sur le pool.

\begin{dated}{2026-09}{Paliers de frais et versions du plus grand AMM}\label{dat:m3:automated-market-makers:uniswap}
Uniswap v2 prélève des frais de 30 points de base sur les transactions, versés aux fournisseurs de liquidité. Uniswap v3 (livre blanc daté de mars 2021, déployé sur
le mainnet d'Ethereum le 5 mai 2021)
introduit la \omterm{def:m3:automated-market-makers:concentrated}{liquidité concentrée} et des paliers de frais, initialement de 0.05~\%, 0.30~\% et 1~\% avec des espacements de 10, 60 et 200 ticks,
où un tick est un pas de prix d'un facteur 1.0001. Uniswap v4 est devenu opérationnel le 31 janvier 2025 et ajoute des hooks, des plugins modulaires
qui permettent aux développeurs d'ajouter une logique sur mesure pour les pools, les swaps, les frais et les positions. Selon le décompte de CoinGecko, le volume au comptant sur les \omterm{def:m3:automated-market-makers:amm}{plateformes décentralisées}
est passé de 6.0~\% du volume au comptant des \omterm{def:m3:centralised-exchanges:cex}{plateformes centralisées} en janvier 2021 à un sommet de 37.4~\% en juin 2025, avant de s'établir près de 21~\% en novembre 2025.
\end{dated}

\section{Liquidité concentrée}

\begin{definition}[Liquidité concentrée]\label{def:m3:automated-market-makers:concentrated}
La \emph{liquidité concentrée}\index{liquidité concentrée} est la liquidité fournie à un \omterm{def:m3:automated-market-makers:amm}{teneur de marché automatisé} sur
une plage de prix choisie uniquement~: la position se comporte comme un \omterm{def:m3:automated-market-makers:cp}{AMM à produit constant} doté d'une réserve virtuelle plus grande tant que le prix est
à l'intérieur de la plage, et ne détient qu'un seul des deux \omterm{def:m3:blockchains-for-traders:key}{jetons} lorsque le prix en sort.
\end{definition}

Dans Uniswap v3, une position est un montant de liquidité $L$ sur une plage $[P_a, P_b)$. À l'intérieur de la plage, avec un prix
en racine carrée $\sqrt{P}$, elle détient $L(1/\sqrt{P} - 1/\sqrt{P_b})$ du premier \omterm{def:m3:blockchains-for-traders:key}{jeton} et $L(\sqrt{P} - \sqrt{P_a})$ du second~; sous
$P_a$, elle ne détient que le premier \omterm{def:m3:blockchains-for-traders:key}{jeton}, au-dessus de $P_b$ seulement le second. Les prix sont stockés sous forme de racines carrées en virgule fixe binaire,
et les plages s'arrêtent sur des ticks. Une plage étroite exige beaucoup moins de \omterm{def:m3:blockchains-for-traders:key}{jetons} qu'une position sur toute la plage pour un même $L$, si bien que le même
capital procure davantage de profondeur près du prix courant~: une plage de $\pm10\%$ est environ 21 fois plus efficace en capital que la
plage complète. Elle vend aussi ses \omterm{def:m3:blockchains-for-traders:key}{jetons} plus vite à mesure que le prix traverse la plage, et ne rapporte rien une fois que le prix
l'a quittée.

\begin{omfigure}
\begin{tikzpicture}
\begin{axis}[width=11cm,height=5cm,xmode=log,log basis x=2,
  xlabel={prix relatif à l'entrée},ylabel={valeur relative à la détention, \%},
  tick label style={font=\small},label style={font=\small},
  xmin=0.5,xmax=2,ymin=-20,ymax=0.5,grid=major,grid style={gray!25},
  xtick={0.5,0.707,1,1.414,2},xticklabels={0.5,0.71,1,1.41,2},
  legend columns=3,legend style={font=\footnotesize,at={(0.5,-0.3)},anchor=north,draw=none,fill=none,/tikz/every even column/.append style={column sep=8pt}},
  table/col sep=comma]
\addplot[omThm,very thick] table[x=r,y=full]{figdata/markets-3/20-automated-market-makers/il.csv};\addlegendentry{plage complète}
\addplot[omDef,very thick,dashed] table[x=r,y=w50]{figdata/markets-3/20-automated-market-makers/il.csv};\addlegendentry{plage $\pm50\%$}
\addplot[omExo,very thick,dotted] table[x=r,y=w10]{figdata/markets-3/20-automated-market-makers/il.csv};\addlegendentry{plage $\pm10\%$}
\end{axis}
\end{tikzpicture}

\omcaption{\omterm{def:m3:automated-market-makers:il}{Perte impermanente}~: valeur d'une position de liquidité comparée à la simple détention des \omterm{def:m3:blockchains-for-traders:key}{jetons} déposés, après
une variation du prix d'un facteur donné, pour une position à produit constant sur toute la plage et pour des positions concentrées sur $\pm50\%$ et $\pm10\%$ autour du prix d'entrée.
Plus la plage est étroite, plus la perte s'accumule vite~; hors de la plage, la position ne détient que le \omterm{def:m3:blockchains-for-traders:key}{jeton} qui a perdu de la valeur
face à l'autre. Données~: le tutoriel du chapitre.}\label{fig:m3:automated-market-makers:il}
\end{omfigure}

\section{Stableswap}

Pour deux \omterm{def:m3:blockchains-for-traders:stable}{stablecoins} censés s'échanger un pour un, le produit constant gaspille la liquidité~: l'essentiel se trouve à des prix
éloignés de un, prix auxquels ils ne devraient jamais s'échanger.

\begin{definition}[Stableswap]\label{def:m3:automated-market-makers:stableswap}
Un \emph{stableswap}\index{stableswap} est un \omterm{def:m3:automated-market-makers:amm}{teneur de marché automatisé} dont l'invariant se comporte comme une somme constante
(sans impact sur le prix) au voisinage de réserves équilibrées et comme un produit constant loin de cet équilibre, combinés par un coefficient d'amplification $A$~:
pour $n$ \omterm{def:m3:blockchains-for-traders:key}{jetons} de soldes $x_i$ et d'invariant $D$,
\[
A n^n \sum_i x_i + D = A D n^n + \frac{D^{n+1}}{n^n \prod_i x_i}.
\]
\end{definition}

L'invariant se résout numériquement, pour $D$ lorsqu'on ajoute de la liquidité et pour un solde lorsqu'une transaction est réalisée, par des
itérations de Newton sur les entiers. Avec $A$ grand, le pool s'échange à peu près un pour un sur une large plage de soldes~; si une pièce subit
un \omterm{def:m3:blockchains-for-traders:stable}{décrochage} (depeg), les arbitragistes la vendent au pool jusqu'à ce que celui-ci ne contienne presque plus que cette pièce, et le prix évolue alors comme dans un pool à produit constant. Les fournisseurs de liquidité
d'un pool de \omterm{def:m3:blockchains-for-traders:stable}{stablecoins} vendent donc un put sur l'ancrage de la pièce la plus faible.

\begin{omfigure}
\begin{tikzpicture}
\begin{axis}[width=11cm,height=5cm,
  xlabel={solde de la pièce X},ylabel={solde de la pièce Y},
  tick label style={font=\small},label style={font=\small},
  xmin=0,xmax=20,ymin=0,ymax=20,grid=major,grid style={gray!25},
  legend columns=3,legend style={font=\footnotesize,at={(0.5,-0.3)},anchor=north,draw=none,fill=none,/tikz/every even column/.append style={column sep=8pt}},
  table/col sep=comma]
\addplot[omThm,very thick] table[x=x,y=cp]{figdata/markets-3/20-automated-market-makers/invariants.csv};\addlegendentry{produit constant}
\addplot[gray,thick,dashed,restrict x to domain=0:10] table[x=x,y=csum]{figdata/markets-3/20-automated-market-makers/invariants.csv};\addlegendentry{somme constante}
\addplot[omDef,very thick,dotted] table[x=x,y=ss]{figdata/markets-3/20-automated-market-makers/invariants.csv};\addlegendentry{stableswap, $A = 85$}
\end{axis}
\end{tikzpicture}

\omcaption{Trois invariants passant par le point d'équilibre (5, 5), d'après la comparaison faite par l'article sur le \omterm{def:m3:automated-market-makers:stableswap}{stableswap}~: le prix est la pente. La
courbe \omterm{def:m3:automated-market-makers:stableswap}{stableswap} suit la somme constante sur une large plage de soldes et ne prend une forme de produit constant que lorsqu'une pièce vient à
manquer. Solveur entier de Curve. Données~: le tutoriel du chapitre.}\label{fig:m3:automated-market-makers:invariants}
\end{omfigure}

\section{Perte impermanente et loss-versus-rebalancing}

\begin{definition}[Perte impermanente]\label{def:m3:automated-market-makers:il}
La \emph{perte impermanente}\index{perte impermanente} d'une position de liquidité est l'écart entre sa valeur et la valeur
des \omterm{def:m3:blockchains-for-traders:key}{jetons} initialement déposés, s'ils avaient simplement été conservés, après une variation de prix~; elle est «~impermanente~» seulement en ce sens qu'elle disparaît si
le prix revient à son niveau de départ.
\end{definition}

\begin{proposition}[Perte impermanente d'une position à produit constant]\label{prop:m3:automated-market-makers:il}
Si le prix du premier \omterm{def:m3:blockchains-for-traders:key}{jeton} varie d'un facteur $r$, une position à produit constant sur toute la plage vaut $2\sqrt{r}/(1 + r)$ fois ce
que vaudraient les \omterm{def:m3:blockchains-for-traders:key}{jetons} initiaux conservés. Une division ou une multiplication par deux du prix coûte 5.7~\%.
\end{proposition}

\begin{proof}
Avec une liquidité $L = \sqrt{xy}$, le pool détient $L/\sqrt{P}$ et $L\sqrt{P}$ et vaut $2L\sqrt{P}$ en second \omterm{def:m3:blockchains-for-traders:key}{jeton}. En partant de
$P_0$ pour aller à $rP_0$, le pool vaut $2L\sqrt{rP_0}$, et les \omterm{def:m3:blockchains-for-traders:key}{jetons} conservés valent $L/\sqrt{P_0} \cdot rP_0 + L\sqrt{P_0} = L\sqrt{P_0}(1 +
r)$. Le rapport vaut $2\sqrt{r}/(1 + r)$.
\end{proof}

La \omterm{def:m3:automated-market-makers:il}{perte impermanente} mêle deux choses~: l'exposition au prix, qu'un hedger pourrait supprimer, et un coût véritable. Milionis, Moallemi,
Roughgarden et Zhang les séparent. Comparons le pool à un portefeuille qui détient, à chaque instant, la même quantité du \omterm{def:m3:blockchains-for-traders:key}{jeton} risqué que le
pool, achetée et vendue aux prix de marché au fil des variations de la détention du pool~: le portefeuille de rééquilibrage porte le même risque de marché que le
pool.

\begin{definition}[Loss-versus-rebalancing]\label{def:m3:automated-market-makers:lvr}
\emph{Loss-versus-rebalancing}\index{loss-versus-rebalancing} (LVR) est l'écart entre la valeur d'un portefeuille de rééquilibrage qui
suit les avoirs d'un \omterm{def:m3:automated-market-makers:amm}{teneur de marché automatisé} en négociant aux prix de marché et la valeur du pool lui-même~; c'est le montant que le pool
perd au profit des arbitragistes en négociant à ses propres prix périmés plutôt qu'à ceux du marché.
\end{definition}

\begin{proposition}[LVR d'un pool à produit constant]\label{prop:m3:automated-market-makers:lvrcp}
Si le prix de marché suit un mouvement brownien géométrique (livre~4) de volatilité $\sigma$ et que les arbitragistes maintiennent un pool à produit constant sans frais
au prix de marché, le LVR s'accumule au taux instantané $\sigma^2/8$ par unité de temps, en fraction de la valeur du
pool. Pour une volatilité quotidienne de 5~\%, cela représente environ 3.1 points de base de la valeur du pool par jour.
\end{proposition}

\begin{proof}
Le résultat général de l'article est un LVR instantané de $\tfrac12\sigma^2P^2\,|x^{*\prime}(P)|$, où $x^*(P)$ est la détention du pool en
\omterm{def:m3:blockchains-for-traders:key}{jeton} risqué. Pour le produit constant, $x^* = L/\sqrt{P}$, $|x^{*\prime}| = \tfrac12 L P^{-3/2}$, et le taux vaut $\tfrac14 L\sigma^2\sqrt{P}$~;
divisé par la valeur du pool $2L\sqrt{P}$, il vaut $\sigma^2/8$.
\end{proof}

Le LVR est la sélection adverse subie par les fournisseurs de liquidité (livre~1, chapitre~19)~: chaque transaction d'arbitrage se fait avec quelqu'un qui sait que le
prix a bougé. Il dépend de la volatilité, et non du volume. Les frais sont ce que le pool gagne auprès de tous les autres. Le tutoriel vérifie la formule fermée par
simulation~: sur 200 trajectoires lognormales à une volatilité annuelle de 60~\%, rééquilibrées toutes les 15 minutes pendant 30 jours, le LVR mesuré atteint en
moyenne 37.1 points de base de la valeur du pool, contre 37.0 pour la formule.

\begin{omfigure}
\begin{tikzpicture}
\begin{axis}[width=11cm,height=5cm,
  xlabel={jour},ylabel={LVR cumulé, pb de valeur du pool},
  tick label style={font=\small},label style={font=\small},
  xmin=0,xmax=30,ymin=0,ymax=40,grid=major,grid style={gray!25},
  legend columns=2,legend style={font=\footnotesize,at={(0.5,-0.3)},anchor=north,draw=none,fill=none,/tikz/every even column/.append style={column sep=8pt}},
  table/col sep=comma]
\addplot[omThm,very thick,mark=*,mark size=1pt] table[x=day,y=measured_bp]{figdata/markets-3/20-automated-market-makers/lvr.csv};\addlegendentry{simulé, moyenne de 200 trajectoires}
\addplot[omDef,thick,dashed] table[x=day,y=theory_bp]{figdata/markets-3/20-automated-market-makers/lvr.csv};\addlegendentry{$\sigma^2/8$ par an}
\end{axis}
\end{tikzpicture}

\omcaption{\omterm{def:m3:automated-market-makers:lvr}{Loss-versus-rebalancing} d'un pool à produit constant sans frais~: le portefeuille de rééquilibrage moins le pool, moyenné sur 200 trajectoires simulées
à une volatilité annuelle de 60~\% avec arbitrage toutes les 15 minutes, comparé à la forme fermée donnée par la \cref{prop:m3:automated-market-makers:lvrcp}~; les deux courbes se confondent. Données~: le tutoriel du
chapitre.}\label{fig:m3:automated-market-makers:lvr}
\end{omfigure}

\begin{proposition}[Quand les frais compensent le LVR]\label{prop:m3:automated-market-makers:breakeven}
Si des traders bruitistes font transiter un volume $V$ par jour dans un pool à produit constant de valeur $W$ avec des frais $\gamma_{\mathrm{fee}}$, les revenus de frais du pool
couvrent son LVR lorsque $V/W \ge \sigma_d^2/(8\gamma_{\mathrm{fee}})$, où $\sigma_d$ est la volatilité quotidienne.
\end{proposition}

\begin{proof}
Les frais perçus valent $\gamma_{\mathrm{fee}}V$ par jour et le LVR $\tfrac18\sigma_d^2 W$ par jour~; il suffit de les comparer.
\end{proof}

À une volatilité annuelle de 60~\%, un pool à 0.30~\% doit faire transiter chaque jour 4.1~\% de sa valeur en volume hors arbitrage pour seulement couvrir le LVR~; un pool à 0.05~\%, 24.7~\%.
\Cref{fig:m3:automated-market-makers:turnover} montre l'exigence en fonction de la volatilité. Concentrer la liquidité augmente les frais perçus et le LVR dans la même
proportion tant que le prix reste dans la plage, de sorte qu'il ne modifie pas le rapport~; cela change en revanche le capital engagé et la fréquence à laquelle la position doit être
déplacée.

\begin{omfigure}
\begin{tikzpicture}
\begin{axis}[width=11cm,height=5cm,ymode=log,log basis y=10,
  xlabel={volatilité annuelle, \%},ylabel={volume quotidien / valeur du pool, \%},
  tick label style={font=\small},label style={font=\small},
  xmin=20,xmax=120,ymin=0.1,ymax=200,grid=major,grid style={gray!25},
  legend columns=3,legend style={font=\footnotesize,at={(0.5,-0.3)},anchor=north,draw=none,fill=none,/tikz/every even column/.append style={column sep=8pt}},
  table/col sep=comma]
\addplot[omThm,very thick] table[x=vol,y=f5]{figdata/markets-3/20-automated-market-makers/turnover.csv};\addlegendentry{frais 0.05~\%}
\addplot[omDef,very thick,dashed] table[x=vol,y=f30]{figdata/markets-3/20-automated-market-makers/turnover.csv};\addlegendentry{frais 0.30~\%}
\addplot[omExo,very thick,dotted] table[x=vol,y=f100]{figdata/markets-3/20-automated-market-makers/turnover.csv};\addlegendentry{frais 1~\%}
\end{axis}
\end{tikzpicture}

\omcaption{Le volume quotidien échangé par les traders de bruit pour lequel les frais d'un pool à produit constant égalent sa \omterm{def:m3:automated-market-makers:lvr}{loss-versus-rebalancing} (LVR)
(\cref{prop:m3:automated-market-makers:breakeven}), pour trois paliers de frais. Données~: le tutoriel du chapitre.}\label{fig:m3:automated-market-makers:turnover}
\end{omfigure}

\section{Liquidité just-in-time, agrégateurs et routage par intentions}

La \omterm{def:m3:automated-market-makers:concentrated}{liquidité concentrée} permet une stratégie que seul l'ordonnancement d'une \omterm{def:m3:blockchains-for-traders:chain}{blockchain} rend possible.

\begin{definition}[Liquidité just-in-time]\label{def:m3:automated-market-makers:jit}
\emph{liquidité just-in-time}\index{liquidité just-in-time} est de la \omterm{def:m3:automated-market-makers:concentrated}{liquidité concentrée} ajoutée dans un intervalle de prix étroit juste avant un gros swap en attente,
dans le même bloc, et retirée juste après celui-ci, de manière à encaisser l'essentiel des frais du swap tout en supportant presque aucune des
pertes d'arbitrage du pool.
\end{definition}

Un fournisseur de \omterm{def:m3:automated-market-makers:jit}{liquidité just-in-time} voit le swap dans le \omterm{def:m3:blockchains-for-traders:mempool}{mempool} (\cref{ch:m3:blockchains-for-traders}) et a besoin que ses propres transactions soient ordonnées avant et
après lui~: c'est une forme des jeux d'ordonnancement du \cref{ch:m3:maximal-extractable-value}. Elle prélève des frais aux fournisseurs
passifs du pool et offre au trader un meilleur prix.

\begin{definition}[Agrégateur de DEX, routage par intentions]\label{def:m3:automated-market-makers:routing}
Un \emph{agrégateur de DEX}\index{agrégateur de DEX} répartit un swap entre plusieurs pools et plateformes de négociation afin de minimiser le coût total pour le trader de l'impact sur le prix,
des frais et du \omterm{def:m3:blockchains-for-traders:gas}{gas}. Le \emph{routage par intentions}\index{routage par intentions} permet au trader de signer un message
indiquant ce qu'il donnera et le minimum qu'il acceptera, et laisse des tiers concurrents (fillers, solveurs) trouver l'exécution et la soumettre sur la \omterm{def:m3:blockchains-for-traders:chain}{blockchain}.
\end{definition}

\begin{dated}{2026-09}{Deux protocoles à base d'intentions}\label{dat:m3:automated-market-makers:intents}
La documentation d'UniswapX décrit des ordres signés qui précisent les montants de sortie d'un swap, exécutés par des fillers concurrents qui utilisent leur propre liquidité ou passent
par d'autres sources~; sur certaines chaînes, l'ordre est une enchère hollandaise dont le prix décroît d'un maximum à un minimum au fil du temps, et le swappeur ne paie
ni \omterm{def:m3:blockchains-for-traders:gas}{gas} ni rien pour les swaps échoués. La documentation de CoW Protocol décrit des intentions signées regroupées en lots, des solveurs tiers se disputant
la meilleure exécution de chaque lot, d'abord en appariant les intentions opposées qui s'y trouvent («~coïncidence des souhaits~») et sinon en passant par
des AMM, des agrégateurs et des teneurs de marché privés.
\end{dated}

Pour une société de trading, les conséquences sont au nombre de deux. En tant que preneur de liquidité, elle peut souvent faire mieux via un agrégateur ou une intention que face à un pool isolé,
et doit les comparer nets des frais de \omterm{def:m3:blockchains-for-traders:gas}{gas}. En tant que fournisseur de liquidité ou filler, elle se dispute avec des sociétés spécialisées le flux d'ordres des intentions, un métier bien plus proche
de la tenue de marché par demande de cotation (RFQ), décrite dans le livre~2, que de la fourniture passive de liquidité en pool.

\section{Tutoriel~: swaps, perte impermanente et LVR}\label{tut:m3:automated-market-makers:lvr}

\begin{tutorial}
\textbf{Objectif.} Implémenter des swaps à produit constant et \omterm{def:m3:automated-market-makers:stableswap}{stableswap} en arithmétique entière, comparer une position de liquidité avec la simple détention des actifs, et mesurer la \omterm{def:m3:automated-market-makers:lvr}{loss-versus-rebalancing}
(LVR) par simulation contre $\sigma^2/8$. \textbf{État final~:} les quatre figures de ce chapitre et les chiffres du problème du week-end.
\begin{tutsteps}
\item \textbf{Le swap.} \Cref{prop:m3:automated-market-makers:swap} en arithmétique entière sur la \omterm{def:m3:blockchains-for-traders:chain}{blockchain}.
\omcode{code/firm/amm/firm_amm.py}{14}{22}{Le swap à produit constant, entrée exacte et sortie exacte.}\label{lst:m3:automated-market-makers:cp}
\item \textbf{LVR.} Détenir la position risquée du pool dans un portefeuille rééquilibré et comparer.
\omcode{code/markets-3/20-automated-market-makers/python/m3_amm.py}{35}{51}{Loss-versus-rebalancing le long d'une trajectoire simulée.}\label{lst:m3:automated-market-makers:lvr}
\item \textbf{Exécuter} \texttt{lvr\_mean()}, \texttt{breakeven\_turnover(0.6, 0.003)} et \texttt{fig\_amm.py}.
\end{tutsteps}
\textbf{Pour aller plus loin.} Doter le pool de frais et ne laisser les arbitragistes trader que lorsque l'écart de prix dépasse ces frais~; mesurer comment la LVR baisse avec le temps de bloc~;
simuler une position concentrée qui est recentrée lorsque le prix sort de son intervalle, et compter ce que coûte le recentrage.
\end{tutorial}

\section{Construction~: la bibliothèque AMM}\label{bld:m3:automated-market-makers:amm}

\begin{build}
\bfield{Rôle} La société miniature doit pricer les swaps exactement comme les contrats les exécuteront (à la dernière unité près), valoriser ses positions de liquidité et
simuler des attaques en sandwich et des liquidations de prêts contre les pools (\cref{ch:m3:maximal-extractable-value,ch:m3:defi-credit-and-leverage}).
\bfield{Interface} \texttt{cp\_amount\_out}, \texttt{cp\_amount\_in}~; \texttt{tick\_at\_price}, \texttt{sqrt\_price\_x96},
\texttt{amounts\_for\_liquidity}, \texttt{next\_sqrt\_price\_from\_amount0}~; \texttt{ss\_get\_d}, \texttt{ss\_get\_y}, \texttt{ss\_amount\_out}. La crate Rust
\texttt{firm\_amm} implémente les fonctions à produit constant et \omterm{def:m3:automated-market-makers:stableswap}{stableswap} en \texttt{u128}.
\bfield{Règles} Entiers uniquement, en unités de base des \omterm{def:m3:blockchains-for-traders:key}{jetons}~; division arrondie vers le bas là où les contrats arrondissent vers le bas, arrondi vers le haut là où ils arrondissent vers le haut~; prix en racine carrée au format Q64.96~;
Les itérations de Newton sont plafonnées et lèvent une exception si elles ne convergent pas.
\bfield{Tests d'acceptation} \texttt{code/firm/amm/tests/} et les tests de la crate Rust~: la formule du swap et la croissance de l'invariant avec les frais~; ticks et montants
de position aux bornes des intervalles~; le déplacement du prix après l'ajout d'un \omterm{def:m3:blockchains-for-traders:key}{jeton}~; l'équilibre du \omterm{def:m3:automated-market-makers:stableswap}{stableswap} et un faible glissement par rapport au produit constant~; un swap avec la même
sortie en Python et en Rust.
\bfield{Pour aller plus loin} Arithmétique exacte des ticks avec des entiers sur 256~bits~; swaps sur plusieurs intervalles franchissant des ticks~; frais et LVR pour une position concentrée recentrée.
\end{build}

\begin{omsources}
\item H.~Adams, N.~Zinsmeister et D.~Robinson, \emph{Uniswap v2 Core}, 2020~; H.~Adams et al., \emph{Uniswap v3 Core}, mars 2021~; Uniswap Labs, «~Uniswap v4 is
here~», 2025.
\item M.~Egorov, \emph{StableSwap --- efficient mechanism for Stablecoin liquidity}, 2019.
\item J.~Milionis, C.~C.~Moallemi, T.~Roughgarden et A.~L.~Zhang, «~Automated Market Making and Loss-Versus-Rebalancing~», arXiv:2208.06046, version de mai
2026.
\item documentation d'UniswapX et de CoW Protocol, septembre 2026~; Uniswap Labs, «~Uniswap v3 Mainnet launch~», 5 mai 2021~; CoinGecko Research, «~DEX to
CEX Volume Ratios~», mis à jour le 17 avril 2026.
\end{omsources}

\section{Exercices}

\begin{exercise}[$\star$]\label{exo:m3:automated-market-makers:1}
Un pool détient 1\,000 ether et 3\,000\,000 dollars avec des frais de 0.30~\%. Que rapporte la vente de 10 ether~? Quel est le prix moyen~?
\end{exercise}

\begin{exercise}[$\star$]\label{exo:m3:automated-market-makers:2}
Quelle est la \omterm{def:m3:automated-market-makers:il}{perte impermanente} d'une position en pleine plage lorsque le prix est multiplié par 4~?
\end{exercise}

\begin{exercise}[$\star$]\label{exo:m3:automated-market-makers:3}
Pourquoi le fournisseur de liquidité d'un pool \omterm{def:m3:automated-market-makers:stableswap}{stableswap} vend-il en pratique un put sur la parité de chaque \omterm{def:m3:blockchains-for-traders:key}{jeton}~?
\end{exercise}

\begin{exercise}[$\star\star$]\label{exo:m3:automated-market-makers:4}
La volatilité annuelle est de 80~\%. Quelle est la LVR d'un pool à produit constant par jour et par an, en fraction de la valeur du pool~?
\end{exercise}

\begin{exercise}[$\star\star$]\label{exo:m3:automated-market-makers:5}
Un pool de 50~millions de dollars avec des frais de 0.05~\% échange 20~millions de dollars de volume hors arbitrage par jour à une volatilité annuelle de 50~\%. Les frais couvrent-ils la LVR~?
\end{exercise}

\begin{exercise}[$\star\star$]\label{exo:m3:automated-market-makers:6}
Expliquer pourquoi la \omterm{def:m3:automated-market-makers:il}{perte impermanente} et la LVR ont la même espérance sous la mesure risque-neutre, mais ne sont pas la même chose.
\end{exercise}

\begin{exercise}[$\star\star\star$]\label{exo:m3:automated-market-makers:7}
\emph{Programmation.} Avec \texttt{lvr\_path}, rééquilibrer toutes les heures au lieu de toutes les 15 minutes. La LVR mesurée change-t-elle~? Qu'en est-il si l'on rééquilibre une fois par jour~?
\end{exercise}

\begin{exercise}[$\star\star\star$]\label{exo:m3:automated-market-makers:8}
\emph{Trouver la faille.} «~Notre pool a rapporté 12~\% par an en frais, donc fournir de la liquidité a rapporté 12~\%.~»
\end{exercise}

\section{Problème~: Faut-il fournir de la liquidité~?}

\begin{problem}[{Problème du week-end --- frais contre loss-versus-rebalancing}]\label{pb:m3:automated-market-makers:1}
Une société de trading envisage de fournir 10~millions de dollars de liquidité à un pool éther--dollar à produit constant, avec des frais de 0.30~\%. La volatilité annuelle de l'éther est de 60~\%.
Le pool, d'une valeur totale de 200~millions de dollars, traite 6~millions de dollars par jour de volume hors arbitrage.

\medskip\noindent\textbf{Partie I --- Les coûts.}
\begin{enumerate}
\item Quelle est la LVR quotidienne du pool, en fraction de sa valeur, et en dollars pour la part de la société~?
\item Quels sont les frais quotidiens de la société~?
\item Quel volume quotidien ferait que les frais compensent la LVR~?
\item Quel est le résultat annuel espéré de la société, en faisant abstraction du risque de marché~?
\item Qu'advient-il de la réponse si la volatilité double~?
\end{enumerate}

\medskip\noindent\textbf{Les alternatives.}
\begin{enumerate}[resume]
\item Le palier à 0.05~\% avec le même volume serait-il meilleur ou moins bon~?
\item Une plage de \omterm{def:m3:automated-market-makers:concentrated}{liquidité concentrée} de $\pm10\%$ modifierait-elle le rapport entre les frais et la LVR~?
\item Que fait la société de son risque de marché, et combien cela coûte-t-il~?
\item Quel volume lui faudrait-il avec des frais de 1~\%~?
\item Comment la \omterm{def:m3:automated-market-makers:jit}{liquidité just-in-time} modifie-t-elle les frais du fournisseur de liquidité passif~?
\end{enumerate}

\medskip\noindent\textbf{Partie III --- Mesure.}
\begin{enumerate}[resume]
\item Comment mesurer le volume hors arbitrage du pool~?
\item Comment mesurer la LVR à partir de l'historique du pool~?
\item Pourquoi le temps de bloc importe-t-il pour la LVR~?
\item Que révèle la \omterm{def:m3:automated-market-makers:il}{perte impermanente} sur un mois au sujet de la LVR~?
\item Que faudrait-il suivre une fois la position ouverte~?
\end{enumerate}

\medskip\noindent\textbf{Partie IV --- Jugement.}
\begin{enumerate}[resume]
\item Qui gagne dans ce pool, et qui perd~?
\item La fourniture passive de liquidité est-elle une activité de teneur de marché~?
\item Pourquoi les pools à fort volume et faible volatilité attirent-ils la liquidité~?
\item Énoncer le \emph{résultat nommé}~: le volume quotidien à partir duquel les frais compensent la LVR, avec une volatilité de 60~\% et des frais de 0.30~\%, et dire si ce pool l'atteint.
\item En une phrase~: que vend un fournisseur de liquidité~?
\end{enumerate}
\end{problem}

\section{Questions d'entretien}

\begin{interviewq}[$\star$ \iqroles{trader}]\label{iq:m3:automated-market-makers:1}
Dériver le montant obtenu lors d'un échange à produit constant avec frais.
\end{interviewq}

\begin{interviewq}[$\star$ \iqroles{chercheur}]\label{iq:m3:automated-market-makers:2}
Qu'est-ce que la \omterm{def:m3:automated-market-makers:il}{perte impermanente}, et pourquoi ce nom est-il trompeur~?
\end{interviewq}

\begin{interviewq}[$\star\star$ \iqroles{chercheur}]\label{iq:m3:automated-market-makers:3}
Dériver la LVR d'un pool à produit constant.
\end{interviewq}

\begin{interviewq}[$\star\star$ \iqroles{développeur}]\label{iq:m3:automated-market-makers:4}
Pourquoi les contrats AMM utilisent-ils l'arithmétique entière et des prix en racine carrée, et que peut-il mal tourner lorsqu'on les réimplémente hors chaîne~?
\end{interviewq}

\begin{interviewq}[$\star\star$ \iqroles{trader}]\label{iq:m3:automated-market-makers:5}
Comment décideriez-vous si un swap de 5~millions de dollars doit passer par un pool, un agrégateur ou un protocole de \omterm{def:m3:automated-market-makers:routing}{routage par intentions}~?
\end{interviewq}

\begin{interviewq}[$\star\star\star$ \iqroles{chercheur, développeur}]\label{iq:m3:automated-market-makers:6}
Concevoir un AMM qui perd moins face aux arbitragistes. Que modifieriez-vous, et que coûterait ce changement~?
\end{interviewq}
